\documentclass[11pt]{article}
\usepackage{amsmath,amssymb,amsthm}
\usepackage[margin=1.1in]{geometry}

\newtheorem{theorem}{Theorem}
\newtheorem{definition}{Definition}
\newtheorem{remark}{Remark}

\title{Laws of Learning: LLN, Concentration, Uniform Convergence, PAC}
\author{Yuval Lavie}
\date{\today}

\begin{document}
\maketitle

\noindent
Why does minimizing error on a \emph{sample} say anything about error on
the \emph{distribution}? This page collects the results that answer that,
in the order they strengthen each other: averages converge (LLN), at a
quantified exponential rate (Hoeffding, Azuma), uniformly over classes
(Glivenko--Cantelli), which is exactly what PAC learnability needs.
\texttt{learning\_theory.py} demonstrates the first three empirically.

\section{Laws of large numbers}

\begin{theorem}[Weak LLN]
If $x_1,\dots,x_n$ are i.i.d.\ with mean $\mu$ and finite variance, then
for every $\varepsilon > 0$,
$\Pr\bigl(|\bar x_n - \mu| > \varepsilon\bigr) \to 0$.
\end{theorem}

\begin{theorem}[Strong LLN]
If the $x_i$ are i.i.d.\ with $\mathbb{E}|x_1| < \infty$, then
$\bar x_n \to \mu$ \emph{almost surely} --- individual sample paths
converge, not just probabilities (visible in the script's first panel).
\end{theorem}

\begin{remark}[Beyond i.i.d.]
Independence can be relaxed: LLNs hold for ergodic stationary sequences
(Birkhoff), for martingale-difference sequences, and under mixing
conditions --- what must survive is that dependence decays fast enough for
averages to stabilize. Azuma's inequality below is the concentration
counterpart for the martingale case, the standard tool once samples are
dependent.
\end{remark}

\section{Concentration: how fast, exactly}

\begin{theorem}[Hoeffding's inequality]
Let $x_1,\dots,x_n$ be independent with $x_i \in [a_i,b_i]$. Then for all
$\varepsilon>0$,
\begin{equation}
  \Pr\Bigl(\,\bigl|\bar x_n - \mathbb{E}\bar x_n\bigr| \ge \varepsilon\Bigr)
  \;\le\;
  2\exp\!\Bigl(-\frac{2n^2\varepsilon^2}{\sum_i (b_i-a_i)^2}\Bigr)
  \;\overset{[0,1]}{=}\;
  2e^{-2n\varepsilon^2}.
  \label{eq:hoeffding}
\end{equation}
\end{theorem}

The LLN says the tail vanishes; Hoeffding says \emph{exponentially fast},
with explicit constants and no variance knowledge --- only boundedness.
Inverting \eqref{eq:hoeffding}: with probability $1-\delta$,
$|\bar x_n - \mu| \le \sqrt{\log(2/\delta)/2n}$ --- the ubiquitous
$1/\sqrt{n}$ of statistics. (Proof route: Chernoff's method,
$\Pr(S\ge\varepsilon) \le e^{-\lambda\varepsilon}\mathbb{E}e^{\lambda S}$,
plus Hoeffding's lemma $\mathbb{E}e^{\lambda(x-\mathbb{E}x)} \le
e^{\lambda^2(b-a)^2/8}$, then optimize $\lambda$.)

\begin{theorem}[Azuma--Hoeffding]
Let $M_0, M_1, \dots$ be a martingale with bounded differences
$|M_k - M_{k-1}| \le c_k$. Then
\begin{equation}
  \Pr\bigl(|M_n - M_0| \ge \varepsilon\bigr)
  \;\le\; 2\exp\!\Bigl(-\frac{\varepsilon^2}{2\sum_k c_k^2}\Bigr).
  \label{eq:azuma}
\end{equation}
\end{theorem}

Hoeffding for dependent-but-martingale increments. Its workhorse corollary
(McDiarmid): any function of independent variables that changes by at most
$c_i$ when one coordinate changes concentrates around its mean --- the
tool behind generalization bounds for quantities that are not plain
averages.

\section{Uniform convergence: Glivenko--Cantelli}

Hoeffding controls \emph{one} fixed quantity. Learning selects a
hypothesis from a whole class \emph{after} seeing data, so we need
convergence \emph{simultaneously} over the class.

\begin{theorem}[Glivenko--Cantelli]
For i.i.d.\ samples with CDF $F$ and empirical CDF
$F_n(t) = \tfrac1n\#\{i : x_i \le t\}$,
\begin{equation}
  \sup_{t\in\mathbb{R}}\, \bigl|F_n(t) - F(t)\bigr|
  \;\xrightarrow{\;a.s.\;}\; 0 ,
  \label{eq:gc}
\end{equation}
and quantitatively (Dvoretzky--Kiefer--Wolfowitz),
$\Pr(\sup_t |F_n - F| > \varepsilon) \le 2e^{-2n\varepsilon^2}$ --- the
same exponent as \eqref{eq:hoeffding}, despite the supremum over
infinitely many events.
\end{theorem}

\begin{definition}[Glivenko--Cantelli class]
A function class $\mathcal{F}$ is GC for a distribution $D$ if
$\sup_{f\in\mathcal{F}} \bigl|\tfrac1n\sum_i f(x_i) -
\mathbb{E}_D f\bigr| \to 0$ (in probability / a.s.).
\end{definition}

The classical theorem says the indicators
$\{\,\mathbf{1}[x\le t] : t\in\mathbb{R}\,\}$ form a GC class. Which
classes are GC is \emph{the} question of learning theory; the answer is
measured by combinatorial capacity --- finite VC dimension implies GC with
rate $\sqrt{(\text{VC}\log n)/n}$ --- and capacity is what
sample-complexity bounds charge for.

\section{PAC learning}

\begin{definition}[PAC learnability]
A hypothesis class $\mathcal{H}$ is PAC learnable if there is an algorithm
$A$ and sample complexity $n(\varepsilon,\delta)$ such that for every
distribution $D$ and every $\varepsilon, \delta \in (0,1)$: given
$n \ge n(\varepsilon,\delta)$ i.i.d.\ samples, with probability at least
$1-\delta$ the output $\hat h = A(S)$ satisfies
\begin{equation}
  L_D(\hat h) \;\le\; \min_{h\in\mathcal{H}} L_D(h) + \varepsilon .
  \label{eq:pac}
\end{equation}
\end{definition}

``Probably ($1-\delta$) Approximately ($\varepsilon$) Correct.'' The
pieces above assemble into the fundamental result: if $\mathcal{H}$ is a
GC class, the \emph{empirical risk minimizer}
$\hat h = \arg\min_{h} \tfrac1n\sum_i \ell(h, x_i)$ PAC-learns
$\mathcal{H}$, because uniform convergence makes empirical risks a
trustworthy proxy for true risks across the whole class at once. Two
standard sample complexities (0--1 loss):
\begin{equation}
  n \;=\; O\!\Bigl(\tfrac{\log|\mathcal{H}| + \log(1/\delta)}
  {\varepsilon^2}\Bigr)
  \;\;\text{(finite class)},
  \qquad
  n \;=\; O\!\Bigl(\tfrac{\mathrm{VC}(\mathcal{H}) + \log(1/\delta)}
  {\varepsilon^2}\Bigr),
  \label{eq:samplecomplexity}
\end{equation}
with $\varepsilon^2 \to \varepsilon$ in the realizable (zero-error) case.
The fundamental theorem of PAC learning closes the loop: for binary
classification, \emph{learnable $\iff$ finite VC dimension $\iff$ uniform
convergence} --- capacity, not cleverness of the algorithm, decides
learnability.

\end{document}
